\section{账本事件的区域制度深描：western-han}

\label{sec:ledger-depth-western-han}

本组叙事以本卷事件账本的可核年序为骨架，将政治节点放回地方资源、制度中介、人口流动与材料类型中。每条来源能够证明的范围不同，故不以朝廷命令、战报数字或后出传记替代基层社会的完整经验。

\subsection{地方资源与制度接口}

账本条目\texttt{WH-0209-DAXIANG}记载，前209年七月，秦末诸起事者发生“陈胜、吴广在大泽乡起事”。这一节点可放在戍卒误期面临秦律惩处，关东对徭役和郡县控制的压力累积的条件下理解；其后续影响被账本概括为反秦动员由局部逃亡转为可复制的政治号召，六国旧地相继响应。就地方尺度而言，事件并不只属于朝廷或统帅：征发、道路、仓储、户籍、市场、寺院与家庭关系都可能决定命令如何抵达、战事如何延长，或接收何以在不同地区呈现不一样的速度。

本条列举的核心材料为《史记·陈涉世家》；《汉书·陈胜项籍传》；《资治通鉴》卷七。这些材料较适合钉住事件的发生、官方表述或特定人物、地点，却不自动构成对全境人口、收入、兵数、信仰或动机的调查。账本的置信与材料提示是“可靠纪年”；来源限度进一步说明“秦末征发、刑律与起事网络的地方社会研究”。因此，不能把条目中的政权名称、行政分类或战报修辞直接转译为固定族群和统一社会意志，也不能将制度宣布写成各地同步实施。

在叙事上，更稳妥的做法是把这一事件视为一条需要地方中介才能运作的链：中央的命令、地方官与精英的选择、运输与劳力的条件、以及普通家庭可能面对的迁徙、赋役或安全风险彼此相连。现有材料能够显示其中若干环节，却保留了大量沉默。承认这种沉默并非放弃解释，而是避免以一条编年记录覆盖不同区域和社群的差异。

账本条目\texttt{WH-0208-XIANG-LIANG}记载，前208年，楚地反秦集团发生“项梁拥立楚怀王后，重建楚国名号”。这一节点可放在陈胜政权瓦解后，反秦武装需要可号召诸地的名义中心的条件下理解；其后续影响被账本概括为楚王名号成为项氏集团动员诸侯和分配军权的重要资源。就地方尺度而言，事件并不只属于朝廷或统帅：征发、道路、仓储、户籍、市场、寺院与家庭关系都可能决定命令如何抵达、战事如何延长，或接收何以在不同地区呈现不一样的速度。

本条列举的核心材料为《史记·项羽本纪》；《史记·陈涉世家》；《汉书·陈胜项籍传》。这些材料较适合钉住事件的发生、官方表述或特定人物、地点，却不自动构成对全境人口、收入、兵数、信仰或动机的调查。账本的置信与材料提示是“王年较可靠；细节有争议”；来源限度进一步说明“楚国记忆、旧贵族与反秦合法性建构研究”。因此，不能把条目中的政权名称、行政分类或战报修辞直接转译为固定族群和统一社会意志，也不能将制度宣布写成各地同步实施。

在叙事上，更稳妥的做法是把这一事件视为一条需要地方中介才能运作的链：中央的命令、地方官与精英的选择、运输与劳力的条件、以及普通家庭可能面对的迁徙、赋役或安全风险彼此相连。现有材料能够显示其中若干环节，却保留了大量沉默。承认这种沉默并非放弃解释，而是避免以一条编年记录覆盖不同区域和社群的差异。

账本条目\texttt{WH-0207-JULU}记载，前207年，楚军与秦军发生“巨鹿之战中项羽军击破秦军主力”。这一节点可放在秦将王离围赵，诸侯军多未敢决战，项羽须以胜利争夺联军领导权的条件下理解；其后续影响被账本概括为秦军关东控制崩解，项羽取得诸侯军事领袖地位。就地方尺度而言，事件并不只属于朝廷或统帅：征发、道路、仓储、户籍、市场、寺院与家庭关系都可能决定命令如何抵达、战事如何延长，或接收何以在不同地区呈现不一样的速度。

本条列举的核心材料为《史记·项羽本纪》；《史记·秦始皇本纪》；《汉书·陈胜项籍传》；《资治通鉴》卷八。这些材料较适合钉住事件的发生、官方表述或特定人物、地点，却不自动构成对全境人口、收入、兵数、信仰或动机的调查。账本的置信与材料提示是“可靠纪年；战果数字有争议”；来源限度进一步说明“巨鹿战役的兵力、补给与诸侯观望研究”。因此，不能把条目中的政权名称、行政分类或战报修辞直接转译为固定族群和统一社会意志，也不能将制度宣布写成各地同步实施。

在叙事上，更稳妥的做法是把这一事件视为一条需要地方中介才能运作的链：中央的命令、地方官与精英的选择、运输与劳力的条件、以及普通家庭可能面对的迁徙、赋役或安全风险彼此相连。现有材料能够显示其中若干环节，却保留了大量沉默。承认这种沉默并非放弃解释，而是避免以一条编年记录覆盖不同区域和社群的差异。

账本条目\texttt{WH-0207-ZHANG-HAN-SURRENDER}记载，前207年冬，秦与楚发生“章邯向项羽降，秦军降卒后遭坑杀的叙事引发争议”。这一节点可放在章邯部在巨鹿失利后补给和政治支持断裂，项羽担忧降军反复的条件下理解；其后续影响被账本概括为关中门户失去野战屏障；降卒遇害的规模与责任成为后世叙事争点。就地方尺度而言，事件并不只属于朝廷或统帅：征发、道路、仓储、户籍、市场、寺院与家庭关系都可能决定命令如何抵达、战事如何延长，或接收何以在不同地区呈现不一样的速度。

本条列举的核心材料为《史记·项羽本纪》；《汉书·陈胜项籍传》；《资治通鉴》卷八。这些材料较适合钉住事件的发生、官方表述或特定人物、地点，却不自动构成对全境人口、收入、兵数、信仰或动机的调查。账本的置信与材料提示是“核心降事实可靠；坑杀规模有争议”；来源限度进一步说明“新安降卒处置、数字夸张与战争暴力记忆研究”。因此，不能把条目中的政权名称、行政分类或战报修辞直接转译为固定族群和统一社会意志，也不能将制度宣布写成各地同步实施。

在叙事上，更稳妥的做法是把这一事件视为一条需要地方中介才能运作的链：中央的命令、地方官与精英的选择、运输与劳力的条件、以及普通家庭可能面对的迁徙、赋役或安全风险彼此相连。现有材料能够显示其中若干环节，却保留了大量沉默。承认这种沉默并非放弃解释，而是避免以一条编年记录覆盖不同区域和社群的差异。

账本条目\texttt{WH-0207-LIU-BANG-GUANZHONG}记载，前207年十月，沛公军与秦发生“刘邦军先入关中，子婴降于轵道”。这一节点可放在秦廷在赵高政变与各路起义夹击下无力守关，刘邦利用武关路线先至的条件下理解；其后续影响被账本概括为秦朝终结，刘邦获得关中行政资源却与项羽发生支配权冲突。就地方尺度而言，事件并不只属于朝廷或统帅：征发、道路、仓储、户籍、市场、寺院与家庭关系都可能决定命令如何抵达、战事如何延长，或接收何以在不同地区呈现不一样的速度。

本条列举的核心材料为《史记·高祖本纪》；《史记·秦始皇本纪》；《汉书·高帝纪》。这些材料较适合钉住事件的发生、官方表述或特定人物、地点，却不自动构成对全境人口、收入、兵数、信仰或动机的调查。账本的置信与材料提示是“可靠纪年；仪式细节有异文”；来源限度进一步说明“秦亡时的地方投降、关中秩序与入关约法研究”。因此，不能把条目中的政权名称、行政分类或战报修辞直接转译为固定族群和统一社会意志，也不能将制度宣布写成各地同步实施。

在叙事上，更稳妥的做法是把这一事件视为一条需要地方中介才能运作的链：中央的命令、地方官与精英的选择、运输与劳力的条件、以及普通家庭可能面对的迁徙、赋役或安全风险彼此相连。现有材料能够显示其中若干环节，却保留了大量沉默。承认这种沉默并非放弃解释，而是避免以一条编年记录覆盖不同区域和社群的差异。

账本条目\texttt{WH-0206-HONGMEN}记载，前206年，项羽集团与刘邦集团发生“鸿门会后项羽控制关中，刘邦暂避其锋”。这一节点可放在刘邦先入关而未能独占关中，项羽以诸侯联军和军事实力压制的条件下理解；其后续影响被账本概括为楚汉两集团的公开决裂暂缓，关中分配转入项羽主导的封建安排。就地方尺度而言，事件并不只属于朝廷或统帅：征发、道路、仓储、户籍、市场、寺院与家庭关系都可能决定命令如何抵达、战事如何延长，或接收何以在不同地区呈现不一样的速度。

本条列举的核心材料为《史记·项羽本纪》；《史记·高祖本纪》；《汉纪》卷一。这些材料较适合钉住事件的发生、官方表述或特定人物、地点，却不自动构成对全境人口、收入、兵数、信仰或动机的调查。账本的置信与材料提示是“事件次序可靠；言辞与细节有争议”；来源限度进一步说明“鸿门叙事的文本层累、范增角色与权力谈判研究”。因此，不能把条目中的政权名称、行政分类或战报修辞直接转译为固定族群和统一社会意志，也不能将制度宣布写成各地同步实施。

在叙事上，更稳妥的做法是把这一事件视为一条需要地方中介才能运作的链：中央的命令、地方官与精英的选择、运输与劳力的条件、以及普通家庭可能面对的迁徙、赋役或安全风险彼此相连。现有材料能够显示其中若干环节，却保留了大量沉默。承认这种沉默并非放弃解释，而是避免以一条编年记录覆盖不同区域和社群的差异。

账本条目\texttt{WH-0206-EIGHTEEN-KINGS}记载，前206年，西楚及诸侯发生“项羽分封十八王，自立西楚霸王，刘邦徙封汉中”。这一节点可放在反秦联军缺少共同的战后分配制度，项羽以军功和亲疏重划诸侯的条件下理解；其后续影响被账本概括为关中被分为三秦，刘邦在汉中积蓄力量；封国间矛盾迅速外显。就地方尺度而言，事件并不只属于朝廷或统帅：征发、道路、仓储、户籍、市场、寺院与家庭关系都可能决定命令如何抵达、战事如何延长，或接收何以在不同地区呈现不一样的速度。

本条列举的核心材料为《史记·项羽本纪》；《汉书·高帝纪》；《汉书·异姓诸侯王表》。这些材料较适合钉住事件的发生、官方表述或特定人物、地点，却不自动构成对全境人口、收入、兵数、信仰或动机的调查。账本的置信与材料提示是“可靠纪年；封国边界须逐项核对”；来源限度进一步说明“楚汉分封的实际控制、封国边界与合法性研究”。因此，不能把条目中的政权名称、行政分类或战报修辞直接转译为固定族群和统一社会意志，也不能将制度宣布写成各地同步实施。

在叙事上，更稳妥的做法是把这一事件视为一条需要地方中介才能运作的链：中央的命令、地方官与精英的选择、运输与劳力的条件、以及普通家庭可能面对的迁徙、赋役或安全风险彼此相连。现有材料能够显示其中若干环节，却保留了大量沉默。承认这种沉默并非放弃解释，而是避免以一条编年记录覆盖不同区域和社群的差异。

账本条目\texttt{WH-0206-THREE-QIN}记载，前206年八月至九月，汉与三秦发生“刘邦出汉中，平定雍、塞、翟三秦”。这一节点可放在项羽将关中分封给秦降将，地方对其统治基础并不稳固的条件下理解；其后续影响被账本概括为刘邦获得关中粮仓和都城区域，形成与西楚长期竞争的财政基础。就地方尺度而言，事件并不只属于朝廷或统帅：征发、道路、仓储、户籍、市场、寺院与家庭关系都可能决定命令如何抵达、战事如何延长，或接收何以在不同地区呈现不一样的速度。

本条列举的核心材料为《史记·高祖本纪》；《史记·淮阴侯列传》；《汉书·高帝纪》。这些材料较适合钉住事件的发生、官方表述或特定人物、地点，却不自动构成对全境人口、收入、兵数、信仰或动机的调查。账本的置信与材料提示是“年代大体可靠；行军路线有争议”；来源限度进一步说明“汉军还定三秦的道路、补给和韩信作用研究”。因此，不能把条目中的政权名称、行政分类或战报修辞直接转译为固定族群和统一社会意志，也不能将制度宣布写成各地同步实施。

在叙事上，更稳妥的做法是把这一事件视为一条需要地方中介才能运作的链：中央的命令、地方官与精英的选择、运输与劳力的条件、以及普通家庭可能面对的迁徙、赋役或安全风险彼此相连。现有材料能够显示其中若干环节，却保留了大量沉默。承认这种沉默并非放弃解释，而是避免以一条编年记录覆盖不同区域和社群的差异。

账本条目\texttt{WH-0205-PENGCHENG}记载，前205年，汉与西楚发生“刘邦联军攻入彭城后遭项羽反击而溃败”。这一节点可放在刘邦扩张过快，诸侯联盟尚未稳固；项羽自齐地回师实施突袭的条件下理解；其后续影响被账本概括为汉军退守荥阳一线，刘邦转向以韩信北进和离间楚盟友的战略。就地方尺度而言，事件并不只属于朝廷或统帅：征发、道路、仓储、户籍、市场、寺院与家庭关系都可能决定命令如何抵达、战事如何延长，或接收何以在不同地区呈现不一样的速度。

本条列举的核心材料为《史记·项羽本纪》；《汉书·高帝纪》；《汉纪》卷一。这些材料较适合钉住事件的发生、官方表述或特定人物、地点，却不自动构成对全境人口、收入、兵数、信仰或动机的调查。账本的置信与材料提示是“可靠纪年；伤亡数为单方叙事”；来源限度进一步说明“彭城战役的联军结构、兵力数字与楚军机动研究”。因此，不能把条目中的政权名称、行政分类或战报修辞直接转译为固定族群和统一社会意志，也不能将制度宣布写成各地同步实施。

在叙事上，更稳妥的做法是把这一事件视为一条需要地方中介才能运作的链：中央的命令、地方官与精英的选择、运输与劳力的条件、以及普通家庭可能面对的迁徙、赋役或安全风险彼此相连。现有材料能够显示其中若干环节，却保留了大量沉默。承认这种沉默并非放弃解释，而是避免以一条编年记录覆盖不同区域和社群的差异。

账本条目\texttt{WH-0205-HEXI-ANNEXATION}记载，前205年，汉与魏发生“韩信、曹参等取魏地，汉军夺取河东交通节点”。这一节点可放在彭城败后汉军需开辟北方战场并切断楚军西北联系的条件下理解；其后续影响被账本概括为汉军取得黄河中游的征发空间，为后续赵、齐战役建立基地。就地方尺度而言，事件并不只属于朝廷或统帅：征发、道路、仓储、户籍、市场、寺院与家庭关系都可能决定命令如何抵达、战事如何延长，或接收何以在不同地区呈现不一样的速度。

本条列举的核心材料为《史记·淮阴侯列传》；《史记·魏豹彭越列传》；《汉书·高帝纪》。这些材料较适合钉住事件的发生、官方表述或特定人物、地点，却不自动构成对全境人口、收入、兵数、信仰或动机的调查。账本的置信与材料提示是“年代与具体将领分工有争议”；来源限度进一步说明“魏地争夺与河东渡口的战略地理研究”。因此，不能把条目中的政权名称、行政分类或战报修辞直接转译为固定族群和统一社会意志，也不能将制度宣布写成各地同步实施。

在叙事上，更稳妥的做法是把这一事件视为一条需要地方中介才能运作的链：中央的命令、地方官与精英的选择、运输与劳力的条件、以及普通家庭可能面对的迁徙、赋役或安全风险彼此相连。现有材料能够显示其中若干环节，却保留了大量沉默。承认这种沉默并非放弃解释，而是避免以一条编年记录覆盖不同区域和社群的差异。

账本条目\texttt{WH-0204-JINGXING}记载，前204年，汉与赵发生“韩信破赵于井陉，收编赵地兵员”。这一节点可放在汉军需在荥阳对峙期间夺取赵地资源，赵军据险防守的条件下理解；其后续影响被账本概括为北方诸国的均势改变，韩信成为汉方独立战区的关键统帅。就地方尺度而言，事件并不只属于朝廷或统帅：征发、道路、仓储、户籍、市场、寺院与家庭关系都可能决定命令如何抵达、战事如何延长，或接收何以在不同地区呈现不一样的速度。

本条列举的核心材料为《史记·淮阴侯列传》；《史记·赵世家》；《汉书·韩信传》。这些材料较适合钉住事件的发生、官方表述或特定人物、地点，却不自动构成对全境人口、收入、兵数、信仰或动机的调查。账本的置信与材料提示是“可靠纪年；背水阵细节有文学化可能”；来源限度进一步说明“井陉战役叙事、太行交通与军事动员研究”。因此，不能把条目中的政权名称、行政分类或战报修辞直接转译为固定族群和统一社会意志，也不能将制度宣布写成各地同步实施。

在叙事上，更稳妥的做法是把这一事件视为一条需要地方中介才能运作的链：中央的命令、地方官与精英的选择、运输与劳力的条件、以及普通家庭可能面对的迁徙、赋役或安全风险彼此相连。现有材料能够显示其中若干环节，却保留了大量沉默。承认这种沉默并非放弃解释，而是避免以一条编年记录覆盖不同区域和社群的差异。

\subsection{道路、边地与人口移动}

账本条目\texttt{WH-0204-XINGYANG}记载，前204年，汉与西楚发生“荥阳、成皋争夺持续，刘邦一度突围西走”。这一节点可放在关中至中原的通道决定双方补给，彭城后汉军主力被楚军牵制的条件下理解；其后续影响被账本概括为楚军未能彻底歼灭刘邦，战争转为多战区的资源竞争。就地方尺度而言，事件并不只属于朝廷或统帅：征发、道路、仓储、户籍、市场、寺院与家庭关系都可能决定命令如何抵达、战事如何延长，或接收何以在不同地区呈现不一样的速度。

本条列举的核心材料为《史记·项羽本纪》；《史记·高祖本纪》；《汉书·高帝纪》；《资治通鉴》卷十。这些材料较适合钉住事件的发生、官方表述或特定人物、地点，却不自动构成对全境人口、收入、兵数、信仰或动机的调查。账本的置信与材料提示是“可靠纪年；纪传对突围人物异同”；来源限度进一步说明“荥阳—成皋走廊、粮道与楚汉持久战研究”。因此，不能把条目中的政权名称、行政分类或战报修辞直接转译为固定族群和统一社会意志，也不能将制度宣布写成各地同步实施。

在叙事上，更稳妥的做法是把这一事件视为一条需要地方中介才能运作的链：中央的命令、地方官与精英的选择、运输与劳力的条件、以及普通家庭可能面对的迁徙、赋役或安全风险彼此相连。现有材料能够显示其中若干环节，却保留了大量沉默。承认这种沉默并非放弃解释，而是避免以一条编年记录覆盖不同区域和社群的差异。

账本条目\texttt{WH-0203-QI-CAMPAIGN}记载，前203年，汉、齐与西楚发生“韩信击齐并在潍水一带破楚援军”。这一节点可放在齐王田广在汉楚之间反复，楚遣龙且援齐以阻断汉军东进的条件下理解；其后续影响被账本概括为汉方获得齐地人口与粮赋，项羽在东部失去可资动员的盟国。就地方尺度而言，事件并不只属于朝廷或统帅：征发、道路、仓储、户籍、市场、寺院与家庭关系都可能决定命令如何抵达、战事如何延长，或接收何以在不同地区呈现不一样的速度。

本条列举的核心材料为《史记·淮阴侯列传》；《史记·田儋列传》；《汉书·高帝纪》。这些材料较适合钉住事件的发生、官方表述或特定人物、地点，却不自动构成对全境人口、收入、兵数、信仰或动机的调查。账本的置信与材料提示是“可靠纪年；水战战术细节有争议”；来源限度进一步说明“齐地征服、潍水战术与田横集团研究”。因此，不能把条目中的政权名称、行政分类或战报修辞直接转译为固定族群和统一社会意志，也不能将制度宣布写成各地同步实施。

在叙事上，更稳妥的做法是把这一事件视为一条需要地方中介才能运作的链：中央的命令、地方官与精英的选择、运输与劳力的条件、以及普通家庭可能面对的迁徙、赋役或安全风险彼此相连。现有材料能够显示其中若干环节，却保留了大量沉默。承认这种沉默并非放弃解释，而是避免以一条编年记录覆盖不同区域和社群的差异。

账本条目\texttt{WH-0203-HONGGOU}记载，前203年，汉与西楚发生“楚汉以鸿沟为界议和，旋即再战”。这一节点可放在楚军长期围攻未获决定性胜利，汉军北方战区已取得优势的条件下理解；其后续影响被账本概括为短暂划界未能化解权力竞争；刘邦联络韩信、彭越准备合围。就地方尺度而言，事件并不只属于朝廷或统帅：征发、道路、仓储、户籍、市场、寺院与家庭关系都可能决定命令如何抵达、战事如何延长，或接收何以在不同地区呈现不一样的速度。

本条列举的核心材料为《史记·项羽本纪》；《史记·高祖本纪》；《资治通鉴》卷十。这些材料较适合钉住事件的发生、官方表述或特定人物、地点，却不自动构成对全境人口、收入、兵数、信仰或动机的调查。账本的置信与材料提示是“核心事实可靠；条款文本不存”；来源限度进一步说明“鸿沟和约的地望、条款与破约责任研究”。因此，不能把条目中的政权名称、行政分类或战报修辞直接转译为固定族群和统一社会意志，也不能将制度宣布写成各地同步实施。

在叙事上，更稳妥的做法是把这一事件视为一条需要地方中介才能运作的链：中央的命令、地方官与精英的选择、运输与劳力的条件、以及普通家庭可能面对的迁徙、赋役或安全风险彼此相连。现有材料能够显示其中若干环节，却保留了大量沉默。承认这种沉默并非放弃解释，而是避免以一条编年记录覆盖不同区域和社群的差异。

账本条目\texttt{WH-0202-GAIXIA}记载，前202年，汉与西楚发生“汉军及诸侯合围项羽于垓下，项羽败死”。这一节点可放在鸿沟后刘邦以齐、梁等兵力合围，楚军补给与盟友相继流失的条件下理解；其后续影响被账本概括为西楚政权瓦解，刘邦成为最有实力的诸侯共主。就地方尺度而言，事件并不只属于朝廷或统帅：征发、道路、仓储、户籍、市场、寺院与家庭关系都可能决定命令如何抵达、战事如何延长，或接收何以在不同地区呈现不一样的速度。

本条列举的核心材料为《史记·项羽本纪》；《汉书·高帝纪》；《资治通鉴》卷十。这些材料较适合钉住事件的发生、官方表述或特定人物、地点，却不自动构成对全境人口、收入、兵数、信仰或动机的调查。账本的置信与材料提示是“可靠纪年；四面楚歌等细节有文学加工”；来源限度进一步说明“垓下战场位置、诸侯参战程度与项羽败亡叙事研究”。因此，不能把条目中的政权名称、行政分类或战报修辞直接转译为固定族群和统一社会意志，也不能将制度宣布写成各地同步实施。

在叙事上，更稳妥的做法是把这一事件视为一条需要地方中介才能运作的链：中央的命令、地方官与精英的选择、运输与劳力的条件、以及普通家庭可能面对的迁徙、赋役或安全风险彼此相连。现有材料能够显示其中若干环节，却保留了大量沉默。承认这种沉默并非放弃解释，而是避免以一条编年记录覆盖不同区域和社群的差异。

账本条目\texttt{WH-0202-HAN-FOUNDATION}记载，前202年二月，汉发生“刘邦即皇帝位，汉朝建立”。这一节点可放在项羽败亡后需要把军事联盟转化为持续的最高统治权的条件下理解；其后续影响被账本概括为汉廷开始配置中央官僚、郡县与异姓诸侯王的复合秩序。就地方尺度而言，事件并不只属于朝廷或统帅：征发、道路、仓储、户籍、市场、寺院与家庭关系都可能决定命令如何抵达、战事如何延长，或接收何以在不同地区呈现不一样的速度。

本条列举的核心材料为《汉书·高帝纪》；《史记·高祖本纪》；《汉纪》卷一。这些材料较适合钉住事件的发生、官方表述或特定人物、地点，却不自动构成对全境人口、收入、兵数、信仰或动机的调查。账本的置信与材料提示是“可靠纪年”；来源限度进一步说明“汉初建国仪式、帝号与战后合法性研究”。因此，不能把条目中的政权名称、行政分类或战报修辞直接转译为固定族群和统一社会意志，也不能将制度宣布写成各地同步实施。

在叙事上，更稳妥的做法是把这一事件视为一条需要地方中介才能运作的链：中央的命令、地方官与精英的选择、运输与劳力的条件、以及普通家庭可能面对的迁徙、赋役或安全风险彼此相连。现有材料能够显示其中若干环节，却保留了大量沉默。承认这种沉默并非放弃解释，而是避免以一条编年记录覆盖不同区域和社群的差异。

账本条目\texttt{WH-0202-CHANGAN}记载，前202年五月，汉发生“汉廷定都长安，依托关中经营全国”。这一节点可放在关中拥有战时粮仓、函谷关险要和秦旧行政网络，洛阳东部位置较难防守的条件下理解；其后续影响被账本概括为长安成为西汉政治中心，关中与关东的资源调配成为长期制度问题。就地方尺度而言，事件并不只属于朝廷或统帅：征发、道路、仓储、户籍、市场、寺院与家庭关系都可能决定命令如何抵达、战事如何延长，或接收何以在不同地区呈现不一样的速度。

本条列举的核心材料为《汉书·高帝纪》；《史记·刘敬叔孙通列传》；汉长安城遗址考古报告。这些材料较适合钉住事件的发生、官方表述或特定人物、地点，却不自动构成对全境人口、收入、兵数、信仰或动机的调查。账本的置信与材料提示是“可靠纪年；城址发展阶段须与考古互证”；来源限度进一步说明“长安建都、关中粮运与帝国空间组织研究”。因此，不能把条目中的政权名称、行政分类或战报修辞直接转译为固定族群和统一社会意志，也不能将制度宣布写成各地同步实施。

在叙事上，更稳妥的做法是把这一事件视为一条需要地方中介才能运作的链：中央的命令、地方官与精英的选择、运输与劳力的条件、以及普通家庭可能面对的迁徙、赋役或安全风险彼此相连。现有材料能够显示其中若干环节，却保留了大量沉默。承认这种沉默并非放弃解释，而是避免以一条编年记录覆盖不同区域和社群的差异。

账本条目\texttt{WH-0201-ENFEOFFMENT}记载，前201年，汉发生“汉廷大规模封建功臣与异姓诸侯王”。这一节点可放在楚汉战争依靠诸侯、将领和地方豪强合作，中央须兑现军功与联盟承诺的条件下理解；其后续影响被账本概括为异姓王拥有军政资源，中央集权与封国安全之间的矛盾制度化。就地方尺度而言，事件并不只属于朝廷或统帅：征发、道路、仓储、户籍、市场、寺院与家庭关系都可能决定命令如何抵达、战事如何延长，或接收何以在不同地区呈现不一样的速度。

本条列举的核心材料为《汉书·高帝纪》；《史记·高祖功臣侯者年表》；《汉书·异姓诸侯王表》。这些材料较适合钉住事件的发生、官方表述或特定人物、地点，却不自动构成对全境人口、收入、兵数、信仰或动机的调查。账本的置信与材料提示是“可靠纪年；封地范围有变动”；来源限度进一步说明“汉初侯国、王国与郡县并置的国家形成研究”。因此，不能把条目中的政权名称、行政分类或战报修辞直接转译为固定族群和统一社会意志，也不能将制度宣布写成各地同步实施。

在叙事上，更稳妥的做法是把这一事件视为一条需要地方中介才能运作的链：中央的命令、地方官与精英的选择、运输与劳力的条件、以及普通家庭可能面对的迁徙、赋役或安全风险彼此相连。现有材料能够显示其中若干环节，却保留了大量沉默。承认这种沉默并非放弃解释，而是避免以一条编年记录覆盖不同区域和社群的差异。

账本条目\texttt{WH-0200-BAIDENG}记载，前200年冬，汉与匈奴发生“刘邦北征受困白登，与匈奴议和撤军”。这一节点可放在韩王信叛汉并引匈奴南下，汉军在北方补给和骑兵条件不足的条件下理解；其后续影响被账本概括为汉廷转向和亲、互市与边郡防御的组合，不再轻易北进。就地方尺度而言，事件并不只属于朝廷或统帅：征发、道路、仓储、户籍、市场、寺院与家庭关系都可能决定命令如何抵达、战事如何延长，或接收何以在不同地区呈现不一样的速度。

本条列举的核心材料为《史记·高祖本纪》；《史记·匈奴列传》；《汉书·高帝纪》；《汉书·匈奴传》。这些材料较适合钉住事件的发生、官方表述或特定人物、地点，却不自动构成对全境人口、收入、兵数、信仰或动机的调查。账本的置信与材料提示是“核心事件可靠；围困天数与和亲条款有争议”；来源限度进一步说明“白登之围、早期和亲与汉匈边境政治研究”。因此，不能把条目中的政权名称、行政分类或战报修辞直接转译为固定族群和统一社会意志，也不能将制度宣布写成各地同步实施。

在叙事上，更稳妥的做法是把这一事件视为一条需要地方中介才能运作的链：中央的命令、地方官与精英的选择、运输与劳力的条件、以及普通家庭可能面对的迁徙、赋役或安全风险彼此相连。现有材料能够显示其中若干环节，却保留了大量沉默。承认这种沉默并非放弃解释，而是避免以一条编年记录覆盖不同区域和社群的差异。

账本条目\texttt{WH-0198-HAN-XIN-REMOVAL}记载，前198年，汉发生“韩信由楚王徙为淮阴侯，异姓王权受抑”。这一节点可放在韩信在楚拥有旧部与声望，刘邦担忧其在齐楚地区的独立动员能力的条件下理解；其后续影响被账本概括为中央以徙封和侯爵化削弱功臣王，为清除异姓王开辟先例。就地方尺度而言，事件并不只属于朝廷或统帅：征发、道路、仓储、户籍、市场、寺院与家庭关系都可能决定命令如何抵达、战事如何延长，或接收何以在不同地区呈现不一样的速度。

本条列举的核心材料为《史记·淮阴侯列传》；《汉书·韩信传》；《资治通鉴》卷十二。这些材料较适合钉住事件的发生、官方表述或特定人物、地点，却不自动构成对全境人口、收入、兵数、信仰或动机的调查。账本的置信与材料提示是“可靠纪年；谋反指控的证据有争议”；来源限度进一步说明“韩信失权、功臣政治与叛乱叙事的史料批判研究”。因此，不能把条目中的政权名称、行政分类或战报修辞直接转译为固定族群和统一社会意志，也不能将制度宣布写成各地同步实施。

在叙事上，更稳妥的做法是把这一事件视为一条需要地方中介才能运作的链：中央的命令、地方官与精英的选择、运输与劳力的条件、以及普通家庭可能面对的迁徙、赋役或安全风险彼此相连。现有材料能够显示其中若干环节，却保留了大量沉默。承认这种沉默并非放弃解释，而是避免以一条编年记录覆盖不同区域和社群的差异。

账本条目\texttt{WH-0196-HAN-XIN-EXECUTION}记载，前196年，汉发生“吕后与萧何诱杀韩信于长乐宫”。这一节点可放在陈豨反叛期间，朝廷将韩信视为潜在的内部威胁的条件下理解；其后续影响被账本概括为异姓功臣集团进一步瓦解，吕后在宫廷决策中的位置上升。就地方尺度而言，事件并不只属于朝廷或统帅：征发、道路、仓储、户籍、市场、寺院与家庭关系都可能决定命令如何抵达、战事如何延长，或接收何以在不同地区呈现不一样的速度。

本条列举的核心材料为《史记·淮阴侯列传》；《汉书·韩信传》；《汉纪》卷三。这些材料较适合钉住事件的发生、官方表述或特定人物、地点，却不自动构成对全境人口、收入、兵数、信仰或动机的调查。账本的置信与材料提示是“可靠纪年；谋反情节有争议”；来源限度进一步说明“韩信案的政治程序、吕后作用与功臣清洗研究”。因此，不能把条目中的政权名称、行政分类或战报修辞直接转译为固定族群和统一社会意志，也不能将制度宣布写成各地同步实施。

在叙事上，更稳妥的做法是把这一事件视为一条需要地方中介才能运作的链：中央的命令、地方官与精英的选择、运输与劳力的条件、以及普通家庭可能面对的迁徙、赋役或安全风险彼此相连。现有材料能够显示其中若干环节，却保留了大量沉默。承认这种沉默并非放弃解释，而是避免以一条编年记录覆盖不同区域和社群的差异。

账本条目\texttt{WH-0196-GUANZHONG-RELOCATION}记载，前196年前后，汉发生“刘敬建议迁徙关东豪族十余万口入关中以实京师”。这一节点可放在关中经历战乱且长安需稳定人力，中央也欲削弱关东豪强的地方根基的条件下理解；其后续影响被账本概括为迁徙成为利用户籍与空间配置巩固京师的手段，其实际覆盖不得外推为全国常态。就地方尺度而言，事件并不只属于朝廷或统帅：征发、道路、仓储、户籍、市场、寺院与家庭关系都可能决定命令如何抵达、战事如何延长，或接收何以在不同地区呈现不一样的速度。

本条列举的核心材料为《史记·刘敬叔孙通列传》；《汉书·高帝纪》；《汉书·地理志》。这些材料较适合钉住事件的发生、官方表述或特定人物、地点，却不自动构成对全境人口、收入、兵数、信仰或动机的调查。账本的置信与材料提示是“人数与执行范围有争议”；来源限度进一步说明“汉初徙民政策、户数口数与关中社会重组研究”。因此，不能把条目中的政权名称、行政分类或战报修辞直接转译为固定族群和统一社会意志，也不能将制度宣布写成各地同步实施。

在叙事上，更稳妥的做法是把这一事件视为一条需要地方中介才能运作的链：中央的命令、地方官与精英的选择、运输与劳力的条件、以及普通家庭可能面对的迁徙、赋役或安全风险彼此相连。现有材料能够显示其中若干环节，却保留了大量沉默。承认这种沉默并非放弃解释，而是避免以一条编年记录覆盖不同区域和社群的差异。

\subsection{社会关系与行政分类}

账本条目\texttt{WH-0196-PENG-YUE}记载，前196年，汉发生“梁王彭越被废杀，梁地改由宗室控制”。这一节点可放在彭越在楚汉战争中掌握梁地兵众，战后中央难以容忍其独立军政能力的条件下理解；其后续影响被账本概括为异姓王封国逐步转入刘氏宗室，中央对东部要地的控制加强。就地方尺度而言，事件并不只属于朝廷或统帅：征发、道路、仓储、户籍、市场、寺院与家庭关系都可能决定命令如何抵达、战事如何延长，或接收何以在不同地区呈现不一样的速度。

本条列举的核心材料为《史记·魏豹彭越列传》；《汉书·彭越传》；《汉书·异姓诸侯王表》。这些材料较适合钉住事件的发生、官方表述或特定人物、地点，却不自动构成对全境人口、收入、兵数、信仰或动机的调查。账本的置信与材料提示是“可靠纪年；罪名与处置叙事有争议”；来源限度进一步说明“彭越案、梁地资源与异姓王处置研究”。因此，不能把条目中的政权名称、行政分类或战报修辞直接转译为固定族群和统一社会意志，也不能将制度宣布写成各地同步实施。

在叙事上，更稳妥的做法是把这一事件视为一条需要地方中介才能运作的链：中央的命令、地方官与精英的选择、运输与劳力的条件、以及普通家庭可能面对的迁徙、赋役或安全风险彼此相连。现有材料能够显示其中若干环节，却保留了大量沉默。承认这种沉默并非放弃解释，而是避免以一条编年记录覆盖不同区域和社群的差异。

账本条目\texttt{WH-0195-YING-BU}记载，前195年，汉与淮南发生“淮南王英布反汉，刘邦亲征后病逝”。这一节点可放在韩信、彭越相继被杀使英布恐惧，淮南又有独立军事资源的条件下理解；其后续影响被账本概括为最后一位强大异姓王败亡；高祖死后皇后与宗室的权力平衡更脆弱。就地方尺度而言，事件并不只属于朝廷或统帅：征发、道路、仓储、户籍、市场、寺院与家庭关系都可能决定命令如何抵达、战事如何延长，或接收何以在不同地区呈现不一样的速度。

本条列举的核心材料为《史记·黥布列传》；《汉书·高帝纪》；《汉书·英布传》。这些材料较适合钉住事件的发生、官方表述或特定人物、地点，却不自动构成对全境人口、收入、兵数、信仰或动机的调查。账本的置信与材料提示是“可靠纪年；作战路线细节有异文”；来源限度进一步说明“英布反叛的联盟崩解、淮南地缘与高祖晚年政治研究”。因此，不能把条目中的政权名称、行政分类或战报修辞直接转译为固定族群和统一社会意志，也不能将制度宣布写成各地同步实施。

在叙事上，更稳妥的做法是把这一事件视为一条需要地方中介才能运作的链：中央的命令、地方官与精英的选择、运输与劳力的条件、以及普通家庭可能面对的迁徙、赋役或安全风险彼此相连。现有材料能够显示其中若干环节，却保留了大量沉默。承认这种沉默并非放弃解释，而是避免以一条编年记录覆盖不同区域和社群的差异。

账本条目\texttt{WH-0195-WHITE-HORSE}记载，前195年前后，汉发生“朝廷以白马盟誓限制非刘氏封王”。这一节点可放在异姓诸侯王叛乱显示封王权力可威胁皇室，刘氏需要排他性合法性语言的条件下理解；其后续影响被账本概括为刘氏宗室封王原则成为后续废立与诛吕行动的重要政治依据。就地方尺度而言，事件并不只属于朝廷或统帅：征发、道路、仓储、户籍、市场、寺院与家庭关系都可能决定命令如何抵达、战事如何延长，或接收何以在不同地区呈现不一样的速度。

本条列举的核心材料为《史记·吕太后本纪》；《汉书·高帝纪下》；《汉书·外戚恩泽侯表》。这些材料较适合钉住事件的发生、官方表述或特定人物、地点，却不自动构成对全境人口、收入、兵数、信仰或动机的调查。账本的置信与材料提示是“盟誓存在可靠；形成过程与效力有争议”；来源限度进一步说明“白马之盟的文本、后世援引与宗室政治研究”。因此，不能把条目中的政权名称、行政分类或战报修辞直接转译为固定族群和统一社会意志，也不能将制度宣布写成各地同步实施。

在叙事上，更稳妥的做法是把这一事件视为一条需要地方中介才能运作的链：中央的命令、地方官与精英的选择、运输与劳力的条件、以及普通家庭可能面对的迁徙、赋役或安全风险彼此相连。现有材料能够显示其中若干环节，却保留了大量沉默。承认这种沉默并非放弃解释，而是避免以一条编年记录覆盖不同区域和社群的差异。

账本条目\texttt{WH-0195-GAOZU-DEATH}记载，前195年四月，汉发生“高祖刘邦去世，惠帝即位，吕后掌握朝政枢纽”。这一节点可放在太子年少而功臣、宗室和皇后家族并存，继承安排未能消除权力竞争的条件下理解；其后续影响被账本概括为吕后通过太后身份介入人事与封爵，汉初政治转入外戚—功臣互动阶段。就地方尺度而言，事件并不只属于朝廷或统帅：征发、道路、仓储、户籍、市场、寺院与家庭关系都可能决定命令如何抵达、战事如何延长，或接收何以在不同地区呈现不一样的速度。

本条列举的核心材料为《汉书·高帝纪下》；《史记·吕太后本纪》；《汉纪》卷三。这些材料较适合钉住事件的发生、官方表述或特定人物、地点，却不自动构成对全境人口、收入、兵数、信仰或动机的调查。账本的置信与材料提示是“可靠纪年”；来源限度进一步说明“高祖死后继承、后宫与开国政治转型研究”。因此，不能把条目中的政权名称、行政分类或战报修辞直接转译为固定族群和统一社会意志，也不能将制度宣布写成各地同步实施。

在叙事上，更稳妥的做法是把这一事件视为一条需要地方中介才能运作的链：中央的命令、地方官与精英的选择、运输与劳力的条件、以及普通家庭可能面对的迁徙、赋役或安全风险彼此相连。现有材料能够显示其中若干环节，却保留了大量沉默。承认这种沉默并非放弃解释，而是避免以一条编年记录覆盖不同区域和社群的差异。

账本条目\texttt{WH-0194-LU-REGIME}记载，前194年至前188年，汉发生“吕后在惠帝朝主导高层人事并处置戚夫人、赵王如意”。这一节点可放在高祖诸子和后宫成员可能成为替代继承中心，吕后力图保护惠帝一支的条件下理解；其后续影响被账本概括为宫廷恐惧与宗室隔阂加深，太后能够直接处分诸王成为既成实践。就地方尺度而言，事件并不只属于朝廷或统帅：征发、道路、仓储、户籍、市场、寺院与家庭关系都可能决定命令如何抵达、战事如何延长，或接收何以在不同地区呈现不一样的速度。

本条列举的核心材料为《史记·吕太后本纪》；《汉书·外戚传上》；《汉纪》卷三。这些材料较适合钉住事件的发生、官方表述或特定人物、地点，却不自动构成对全境人口、收入、兵数、信仰或动机的调查。账本的置信与材料提示是“核心处置事实可靠；叙事细节与动机有争议”；来源限度进一步说明“吕后政治、宫廷暴力叙事与母后权力研究”。因此，不能把条目中的政权名称、行政分类或战报修辞直接转译为固定族群和统一社会意志，也不能将制度宣布写成各地同步实施。

在叙事上，更稳妥的做法是把这一事件视为一条需要地方中介才能运作的链：中央的命令、地方官与精英的选择、运输与劳力的条件、以及普通家庭可能面对的迁徙、赋役或安全风险彼此相连。现有材料能够显示其中若干环节，却保留了大量沉默。承认这种沉默并非放弃解释，而是避免以一条编年记录覆盖不同区域和社群的差异。

账本条目\texttt{WH-0188-HUIDI-DEATH}记载，前188年八月，汉发生“惠帝去世，吕后临朝并先后立少帝”。这一节点可放在惠帝无能独立统御诸侯与功臣，吕后已掌握宫廷与禁军网络的条件下理解；其后续影响被账本概括为皇帝权威被太后临朝结构重塑，刘氏宗室对吕氏扩权的疑虑加深。就地方尺度而言，事件并不只属于朝廷或统帅：征发、道路、仓储、户籍、市场、寺院与家庭关系都可能决定命令如何抵达、战事如何延长，或接收何以在不同地区呈现不一样的速度。

本条列举的核心材料为《史记·吕太后本纪》；《汉书·高后纪》；《汉纪》卷四。这些材料较适合钉住事件的发生、官方表述或特定人物、地点，却不自动构成对全境人口、收入、兵数、信仰或动机的调查。账本的置信与材料提示是“可靠纪年；少帝身份和年龄有争议”；来源限度进一步说明“惠帝死后少帝身份、太后临朝与继承程序研究”。因此，不能把条目中的政权名称、行政分类或战报修辞直接转译为固定族群和统一社会意志，也不能将制度宣布写成各地同步实施。

在叙事上，更稳妥的做法是把这一事件视为一条需要地方中介才能运作的链：中央的命令、地方官与精英的选择、运输与劳力的条件、以及普通家庭可能面对的迁徙、赋役或安全风险彼此相连。现有材料能够显示其中若干环节，却保留了大量沉默。承认这种沉默并非放弃解释，而是避免以一条编年记录覆盖不同区域和社群的差异。

账本条目\texttt{WH-0184-LU-KINGS}记载，前184年前后，汉发生“吕后封吕氏诸王并任用吕产、吕禄等掌军政”。这一节点可放在吕后须为少帝政权建立可靠亲族支柱，同时面对刘氏封王的制约的条件下理解；其后续影响被账本概括为吕氏进入王国与中央军权，太后死后爆发政权归属危机的条件形成。就地方尺度而言，事件并不只属于朝廷或统帅：征发、道路、仓储、户籍、市场、寺院与家庭关系都可能决定命令如何抵达、战事如何延长，或接收何以在不同地区呈现不一样的速度。

本条列举的核心材料为《史记·吕太后本纪》；《汉书·高后纪》；《汉书·外戚传上》。这些材料较适合钉住事件的发生、官方表述或特定人物、地点，却不自动构成对全境人口、收入、兵数、信仰或动机的调查。账本的置信与材料提示是“相对次序可靠；封王年份有异说”；来源限度进一步说明“吕氏封王、军权配置与白马盟政治研究”。因此，不能把条目中的政权名称、行政分类或战报修辞直接转译为固定族群和统一社会意志，也不能将制度宣布写成各地同步实施。

在叙事上，更稳妥的做法是把这一事件视为一条需要地方中介才能运作的链：中央的命令、地方官与精英的选择、运输与劳力的条件、以及普通家庭可能面对的迁徙、赋役或安全风险彼此相连。现有材料能够显示其中若干环节，却保留了大量沉默。承认这种沉默并非放弃解释，而是避免以一条编年记录覆盖不同区域和社群的差异。

账本条目\texttt{WH-0180-ZHULU}记载，前180年八月，汉发生“吕后去世后，周勃、陈平等发动政变诛灭吕氏”。这一节点可放在吕氏依赖太后个人权威，少帝与吕氏王侯难获广泛承认；功臣援引刘氏盟誓的条件下理解；其后续影响被账本概括为吕氏政权瓦解，功臣集团转而选择代王刘恒继位。就地方尺度而言，事件并不只属于朝廷或统帅：征发、道路、仓储、户籍、市场、寺院与家庭关系都可能决定命令如何抵达、战事如何延长，或接收何以在不同地区呈现不一样的速度。

本条列举的核心材料为《史记·吕太后本纪》；《汉书·高后纪》；《汉书·陈丞相世家》。这些材料较适合钉住事件的发生、官方表述或特定人物、地点，却不自动构成对全境人口、收入、兵数、信仰或动机的调查。账本的置信与材料提示是“可靠纪年；参与者分工与宫廷过程有争议”；来源限度进一步说明“诛吕政变、北军南军与功臣—宗室联盟研究”。因此，不能把条目中的政权名称、行政分类或战报修辞直接转译为固定族群和统一社会意志，也不能将制度宣布写成各地同步实施。

在叙事上，更稳妥的做法是把这一事件视为一条需要地方中介才能运作的链：中央的命令、地方官与精英的选择、运输与劳力的条件、以及普通家庭可能面对的迁徙、赋役或安全风险彼此相连。现有材料能够显示其中若干环节，却保留了大量沉默。承认这种沉默并非放弃解释，而是避免以一条编年记录覆盖不同区域和社群的差异。

账本条目\texttt{WH-0180-WENDI-ACCESSION}记载，前180年九月，汉发生“代王刘恒入长安即帝位，是为文帝”。这一节点可放在诛吕者需一位有刘氏血统而外家势力较弱的成年皇帝，以安定宗室和功臣的条件下理解；其后续影响被账本概括为代国王进入中央，西汉以宗室继承修复政权合法性。就地方尺度而言，事件并不只属于朝廷或统帅：征发、道路、仓储、户籍、市场、寺院与家庭关系都可能决定命令如何抵达、战事如何延长，或接收何以在不同地区呈现不一样的速度。

本条列举的核心材料为《汉书·文帝纪》；《史记·孝文本纪》；《汉纪》卷五。这些材料较适合钉住事件的发生、官方表述或特定人物、地点，却不自动构成对全境人口、收入、兵数、信仰或动机的调查。账本的置信与材料提示是“可靠纪年”；来源限度进一步说明“文帝选择、代国资源与诛吕后权力再平衡研究”。因此，不能把条目中的政权名称、行政分类或战报修辞直接转译为固定族群和统一社会意志，也不能将制度宣布写成各地同步实施。

在叙事上，更稳妥的做法是把这一事件视为一条需要地方中介才能运作的链：中央的命令、地方官与精英的选择、运输与劳力的条件、以及普通家庭可能面对的迁徙、赋役或安全风险彼此相连。现有材料能够显示其中若干环节，却保留了大量沉默。承认这种沉默并非放弃解释，而是避免以一条编年记录覆盖不同区域和社群的差异。

账本条目\texttt{WH-0178-LAND-TAX}记载，前178年，汉发生“文帝诏减田租为三十税一”。这一节点可放在战后人口、土地和生产仍在恢复，中央需减少农户负担并稳定编户的条件下理解；其后续影响被账本概括为轻徭薄赋成为文景政治记忆的一部分；各地实际税率与附加役使仍需文书检验。就地方尺度而言，事件并不只属于朝廷或统帅：征发、道路、仓储、户籍、市场、寺院与家庭关系都可能决定命令如何抵达、战事如何延长，或接收何以在不同地区呈现不一样的速度。

本条列举的核心材料为《汉书·文帝纪》；《汉书·食货志》；《汉纪》卷六。这些材料较适合钉住事件的发生、官方表述或特定人物、地点，却不自动构成对全境人口、收入、兵数、信仰或动机的调查。账本的置信与材料提示是“可靠纪年；实际征收程度有争议”；来源限度进一步说明“文帝减租的诏令、地方执行与汉初财政能力研究”。因此，不能把条目中的政权名称、行政分类或战报修辞直接转译为固定族群和统一社会意志，也不能将制度宣布写成各地同步实施。

在叙事上，更稳妥的做法是把这一事件视为一条需要地方中介才能运作的链：中央的命令、地方官与精英的选择、运输与劳力的条件、以及普通家庭可能面对的迁徙、赋役或安全风险彼此相连。现有材料能够显示其中若干环节，却保留了大量沉默。承认这种沉默并非放弃解释，而是避免以一条编年记录覆盖不同区域和社群的差异。

账本条目\texttt{WH-0177-PRIVATE-COINAGE}记载，前177年前后，汉发生“文帝弛山泽禁并许可民间铸钱，后引发货币质量问题”。这一节点可放在中央财政有限且希望减轻管制，文帝朝调整秦以来的禁令的条件下理解；其后续影响被账本概括为民间与诸侯铸钱扩大，币制质量和地方财力问题为武帝收铸埋下伏笔。就地方尺度而言，事件并不只属于朝廷或统帅：征发、道路、仓储、户籍、市场、寺院与家庭关系都可能决定命令如何抵达、战事如何延长，或接收何以在不同地区呈现不一样的速度。

本条列举的核心材料为《汉书·食货志》；《史记·平准书》；西汉钱币范铸遗存。这些材料较适合钉住事件的发生、官方表述或特定人物、地点，却不自动构成对全境人口、收入、兵数、信仰或动机的调查。账本的置信与材料提示是“年代大体可靠；政策范围有争议”；来源限度进一步说明“汉初私铸、铜料供应与货币多样化研究”。因此，不能把条目中的政权名称、行政分类或战报修辞直接转译为固定族群和统一社会意志，也不能将制度宣布写成各地同步实施。

在叙事上，更稳妥的做法是把这一事件视为一条需要地方中介才能运作的链：中央的命令、地方官与精英的选择、运输与劳力的条件、以及普通家庭可能面对的迁徙、赋役或安全风险彼此相连。现有材料能够显示其中若干环节，却保留了大量沉默。承认这种沉默并非放弃解释，而是避免以一条编年记录覆盖不同区域和社群的差异。

\subsection{信仰、知识与物质空间}

账本条目\texttt{WH-0167-CORPORAL-PUNISHMENT}记载，前167年，汉发生“文帝废除肉刑，改以笞刑等处罚”。这一节点可放在缇萦上书故事所反映的刑罚讨论与统治者减轻酷刑的政治需要相互交织的条件下理解；其后续影响被账本概括为法定肉刑被废，但笞刑死亡和地方司法实践显示刑罚负担并未简单消失。就地方尺度而言，事件并不只属于朝廷或统帅：征发、道路、仓储、户籍、市场、寺院与家庭关系都可能决定命令如何抵达、战事如何延长，或接收何以在不同地区呈现不一样的速度。

本条列举的核心材料为《汉书·文帝纪》；《汉书·刑法志》；《汉纪》卷八。这些材料较适合钉住事件的发生、官方表述或特定人物、地点，却不自动构成对全境人口、收入、兵数、信仰或动机的调查。账本的置信与材料提示是“可靠纪年；实际司法效果有争议”；来源限度进一步说明“汉文帝刑制改革、笞数与司法执行研究”。因此，不能把条目中的政权名称、行政分类或战报修辞直接转译为固定族群和统一社会意志，也不能将制度宣布写成各地同步实施。

在叙事上，更稳妥的做法是把这一事件视为一条需要地方中介才能运作的链：中央的命令、地方官与精英的选择、运输与劳力的条件、以及普通家庭可能面对的迁徙、赋役或安全风险彼此相连。现有材料能够显示其中若干环节，却保留了大量沉默。承认这种沉默并非放弃解释，而是避免以一条编年记录覆盖不同区域和社群的差异。

账本条目\texttt{WH-0166-XIONGNU-INVASION}记载，前166年，汉与匈奴发生“匈奴入寇至雍地，文帝调兵守长安西北”。这一节点可放在和亲未消除边境竞争，匈奴可利用草原机动绕过边郡防线的条件下理解；其后续影响被账本概括为长安安全直接暴露，朝廷加强边郡驻军和将领选任的讨论。就地方尺度而言，事件并不只属于朝廷或统帅：征发、道路、仓储、户籍、市场、寺院与家庭关系都可能决定命令如何抵达、战事如何延长，或接收何以在不同地区呈现不一样的速度。

本条列举的核心材料为《汉书·文帝纪》；《汉书·匈奴传》；《史记·匈奴列传》。这些材料较适合钉住事件的发生、官方表述或特定人物、地点，却不自动构成对全境人口、收入、兵数、信仰或动机的调查。账本的置信与材料提示是“可靠纪年；入寇兵数有争议”；来源限度进一步说明“文帝时期汉匈边防、骑兵与和亲体制研究”。因此，不能把条目中的政权名称、行政分类或战报修辞直接转译为固定族群和统一社会意志，也不能将制度宣布写成各地同步实施。

在叙事上，更稳妥的做法是把这一事件视为一条需要地方中介才能运作的链：中央的命令、地方官与精英的选择、运输与劳力的条件、以及普通家庭可能面对的迁徙、赋役或安全风险彼此相连。现有材料能够显示其中若干环节，却保留了大量沉默。承认这种沉默并非放弃解释，而是避免以一条编年记录覆盖不同区域和社群的差异。

账本条目\texttt{WH-0165-JIA-YI}记载，前165年前后，汉发生“贾谊上论削藩与诸侯王问题，未获全面采纳”。这一节点可放在同姓诸侯国拥有广阔领地和独立官属，中央已担忧其尾大不掉的条件下理解；其后续影响被账本概括为削藩成为朝廷可讨论的政策选项，但文帝朝未以强制手段全面实施。就地方尺度而言，事件并不只属于朝廷或统帅：征发、道路、仓储、户籍、市场、寺院与家庭关系都可能决定命令如何抵达、战事如何延长，或接收何以在不同地区呈现不一样的速度。

本条列举的核心材料为《汉书·贾谊传》；《史记·屈原贾生列传》；《汉纪》卷八。这些材料较适合钉住事件的发生、官方表述或特定人物、地点，却不自动构成对全境人口、收入、兵数、信仰或动机的调查。账本的置信与材料提示是“年代与文本成篇过程有争议”；来源限度进一步说明“贾谊政论的成篇年代、削藩方案与文帝政治研究”。因此，不能把条目中的政权名称、行政分类或战报修辞直接转译为固定族群和统一社会意志，也不能将制度宣布写成各地同步实施。

在叙事上，更稳妥的做法是把这一事件视为一条需要地方中介才能运作的链：中央的命令、地方官与精英的选择、运输与劳力的条件、以及普通家庭可能面对的迁徙、赋役或安全风险彼此相连。现有材料能够显示其中若干环节，却保留了大量沉默。承认这种沉默并非放弃解释，而是避免以一条编年记录覆盖不同区域和社群的差异。

账本条目\texttt{WH-0162-ROYAL-PRINCES}记载，前162年前后，汉发生“文帝分封诸子为王，宗室王国网络扩大”。这一节点可放在刘氏政权需要以宗室控制要地，也需安置皇子和维系血统联盟的条件下理解；其后续影响被账本概括为同姓封国数量增加，中央与王国之间的财政、司法和继承摩擦继续累积。就地方尺度而言，事件并不只属于朝廷或统帅：征发、道路、仓储、户籍、市场、寺院与家庭关系都可能决定命令如何抵达、战事如何延长，或接收何以在不同地区呈现不一样的速度。

本条列举的核心材料为《汉书·诸侯王表》；《汉书·文帝纪》；《史记·孝文本纪》。这些材料较适合钉住事件的发生、官方表述或特定人物、地点，却不自动构成对全境人口、收入、兵数、信仰或动机的调查。账本的置信与材料提示是“相对次序可靠；各王受封年有差异”；来源限度进一步说明“文帝宗室分封、王国治理与继承风险研究”。因此，不能把条目中的政权名称、行政分类或战报修辞直接转译为固定族群和统一社会意志，也不能将制度宣布写成各地同步实施。

在叙事上，更稳妥的做法是把这一事件视为一条需要地方中介才能运作的链：中央的命令、地方官与精英的选择、运输与劳力的条件、以及普通家庭可能面对的迁徙、赋役或安全风险彼此相连。现有材料能够显示其中若干环节，却保留了大量沉默。承认这种沉默并非放弃解释，而是避免以一条编年记录覆盖不同区域和社群的差异。

账本条目\texttt{WH-0157-JINGDI-ACCESSION}记载，前157年六月，汉发生“文帝去世，太子启即位，是为景帝”。这一节点可放在文帝以太子启为稳定继承人，功臣政治已转向以宗室皇帝为中心的条件下理解；其后续影响被账本概括为景帝承接减税与边防政策，也面对诸侯王问题的紧迫化。就地方尺度而言，事件并不只属于朝廷或统帅：征发、道路、仓储、户籍、市场、寺院与家庭关系都可能决定命令如何抵达、战事如何延长，或接收何以在不同地区呈现不一样的速度。

本条列举的核心材料为《汉书·景帝纪》；《史记·孝景本纪》；《汉纪》卷九。这些材料较适合钉住事件的发生、官方表述或特定人物、地点，却不自动构成对全境人口、收入、兵数、信仰或动机的调查。账本的置信与材料提示是“可靠纪年”；来源限度进一步说明“文景继承、太子教育与宫廷连续性研究”。因此，不能把条目中的政权名称、行政分类或战报修辞直接转译为固定族群和统一社会意志，也不能将制度宣布写成各地同步实施。

在叙事上，更稳妥的做法是把这一事件视为一条需要地方中介才能运作的链：中央的命令、地方官与精英的选择、运输与劳力的条件、以及普通家庭可能面对的迁徙、赋役或安全风险彼此相连。现有材料能够显示其中若干环节，却保留了大量沉默。承认这种沉默并非放弃解释，而是避免以一条编年记录覆盖不同区域和社群的差异。

账本条目\texttt{WH-0154-CUOFAN}记载，前154年春，汉与吴楚诸王发生“御史大夫晁错推行削藩，吴楚等七国举兵”。这一节点可放在中央以违法为由削减王国封地，吴王刘濞等将既有自治危机转化为武装反抗的条件下理解；其后续影响被账本概括为诸侯王与中央的矛盾由制度争议升级为全国性内战。就地方尺度而言，事件并不只属于朝廷或统帅：征发、道路、仓储、户籍、市场、寺院与家庭关系都可能决定命令如何抵达、战事如何延长，或接收何以在不同地区呈现不一样的速度。

本条列举的核心材料为《汉书·景帝纪》；《汉书·晁错传》；《史记·吴王濞列传》。这些材料较适合钉住事件的发生、官方表述或特定人物、地点，却不自动构成对全境人口、收入、兵数、信仰或动机的调查。账本的置信与材料提示是“可靠纪年；削地程序与先后有争议”；来源限度进一步说明“七国之乱的削地政策、王国财政与叛乱联盟研究”。因此，不能把条目中的政权名称、行政分类或战报修辞直接转译为固定族群和统一社会意志，也不能将制度宣布写成各地同步实施。

在叙事上，更稳妥的做法是把这一事件视为一条需要地方中介才能运作的链：中央的命令、地方官与精英的选择、运输与劳力的条件、以及普通家庭可能面对的迁徙、赋役或安全风险彼此相连。现有材料能够显示其中若干环节，却保留了大量沉默。承认这种沉默并非放弃解释，而是避免以一条编年记录覆盖不同区域和社群的差异。

账本条目\texttt{WH-0154-QIZHIGUO}记载，前154年夏，汉与吴楚七国发生“周亚夫守昌邑、断吴楚粮道，七国之乱被平定”。这一节点可放在吴楚军急于西进，汉廷以梁国牵制并调集中央和诸侯兵力的条件下理解；其后续影响被账本概括为中央避免正面决战而瓦解叛军，王国军政权限随后被持续削弱。就地方尺度而言，事件并不只属于朝廷或统帅：征发、道路、仓储、户籍、市场、寺院与家庭关系都可能决定命令如何抵达、战事如何延长，或接收何以在不同地区呈现不一样的速度。

本条列举的核心材料为《汉书·景帝纪》；《史记·绛侯周勃世家》；《汉书·周亚夫传》。这些材料较适合钉住事件的发生、官方表述或特定人物、地点，却不自动构成对全境人口、收入、兵数、信仰或动机的调查。账本的置信与材料提示是“可靠纪年；战役细节有异文”；来源限度进一步说明“周亚夫战略、昌邑位置与吴楚粮道研究”。因此，不能把条目中的政权名称、行政分类或战报修辞直接转译为固定族群和统一社会意志，也不能将制度宣布写成各地同步实施。

在叙事上，更稳妥的做法是把这一事件视为一条需要地方中介才能运作的链：中央的命令、地方官与精英的选择、运输与劳力的条件、以及普通家庭可能面对的迁徙、赋役或安全风险彼此相连。现有材料能够显示其中若干环节，却保留了大量沉默。承认这种沉默并非放弃解释，而是避免以一条编年记录覆盖不同区域和社群的差异。

账本条目\texttt{WH-0154-CHAO-CUO}记载，前154年，汉发生“景帝为安抚叛王而诛晁错，仍未阻止叛军”。这一节点可放在朝廷误判诛杀倡议者可使吴楚退兵，叛军实际目标已超出政策撤回的条件下理解；其后续影响被账本概括为削藩争议转为不可逆的军事冲突，也显示皇帝对大臣的责任转移。就地方尺度而言，事件并不只属于朝廷或统帅：征发、道路、仓储、户籍、市场、寺院与家庭关系都可能决定命令如何抵达、战事如何延长，或接收何以在不同地区呈现不一样的速度。

本条列举的核心材料为《汉书·晁错传》；《汉书·景帝纪》；《资治通鉴》卷十六。这些材料较适合钉住事件的发生、官方表述或特定人物、地点，却不自动构成对全境人口、收入、兵数、信仰或动机的调查。账本的置信与材料提示是“可靠纪年；诏令动机有争议”；来源限度进一步说明“晁错被杀的决策信息、袁盎建议与危机政治研究”。因此，不能把条目中的政权名称、行政分类或战报修辞直接转译为固定族群和统一社会意志，也不能将制度宣布写成各地同步实施。

在叙事上，更稳妥的做法是把这一事件视为一条需要地方中介才能运作的链：中央的命令、地方官与精英的选择、运输与劳力的条件、以及普通家庭可能面对的迁徙、赋役或安全风险彼此相连。现有材料能够显示其中若干环节，却保留了大量沉默。承认这种沉默并非放弃解释，而是避免以一条编年记录覆盖不同区域和社群的差异。

账本条目\texttt{WH-0144-PRINCELY-REFORM}记载，前144年前后，汉发生“七国之乱后，王国官吏任免和司法权进一步归中央节制”。这一节点可放在叛乱表明王国可依靠自身官属、赋役和兵力挑战中央的条件下理解；其后续影响被账本概括为同姓王保留封号和食邑而军政自主空间收缩，郡国并行转为更不对称的结构。就地方尺度而言，事件并不只属于朝廷或统帅：征发、道路、仓储、户籍、市场、寺院与家庭关系都可能决定命令如何抵达、战事如何延长，或接收何以在不同地区呈现不一样的速度。

本条列举的核心材料为《汉书·百官公卿表》；《汉书·诸侯王表》；《汉书·景帝纪》。这些材料较适合钉住事件的发生、官方表述或特定人物、地点，却不自动构成对全境人口、收入、兵数、信仰或动机的调查。账本的置信与材料提示是“制度演变可靠；具体年份和地域有争议”；来源限度进一步说明“汉代王国官制、内史体系与中央化程度研究”。因此，不能把条目中的政权名称、行政分类或战报修辞直接转译为固定族群和统一社会意志，也不能将制度宣布写成各地同步实施。

在叙事上，更稳妥的做法是把这一事件视为一条需要地方中介才能运作的链：中央的命令、地方官与精英的选择、运输与劳力的条件、以及普通家庭可能面对的迁徙、赋役或安全风险彼此相连。现有材料能够显示其中若干环节，却保留了大量沉默。承认这种沉默并非放弃解释，而是避免以一条编年记录覆盖不同区域和社群的差异。

账本条目\texttt{WH-0141-WUDI-ACCESSION}记载，前141年正月，汉发生“景帝去世，太子彻即位，是为武帝”。这一节点可放在景帝完成削藩的初步框架，年轻皇帝可在较强中央财政和军权基础上扩张的条件下理解；其后续影响被账本概括为朝廷人事和意识形态争论加剧，武帝时代的对外战争与制度重组展开。就地方尺度而言，事件并不只属于朝廷或统帅：征发、道路、仓储、户籍、市场、寺院与家庭关系都可能决定命令如何抵达、战事如何延长，或接收何以在不同地区呈现不一样的速度。

本条列举的核心材料为《汉书·武帝纪》；《史记·孝武本纪》；《汉纪》卷十三。这些材料较适合钉住事件的发生、官方表述或特定人物、地点，却不自动构成对全境人口、收入、兵数、信仰或动机的调查。账本的置信与材料提示是“可靠纪年”；来源限度进一步说明“武帝继位、外戚与儒术政策转向研究”。因此，不能把条目中的政权名称、行政分类或战报修辞直接转译为固定族群和统一社会意志，也不能将制度宣布写成各地同步实施。

在叙事上，更稳妥的做法是把这一事件视为一条需要地方中介才能运作的链：中央的命令、地方官与精英的选择、运输与劳力的条件、以及普通家庭可能面对的迁徙、赋役或安全风险彼此相连。现有材料能够显示其中若干环节，却保留了大量沉默。承认这种沉默并非放弃解释，而是避免以一条编年记录覆盖不同区域和社群的差异。

账本条目\texttt{WH-0140-JIANYUAN}记载，前140年，汉发生“武帝改元建元，任用赵绾、王臧并一度改革礼制议程”。这一节点可放在新帝希望借经学人士调整朝廷议程，窦太后仍掌握深厚的政治影响的条件下理解；其后续影响被账本概括为早期改革受挫，却使经学、人事与皇权关系成为武帝朝持续议题。就地方尺度而言，事件并不只属于朝廷或统帅：征发、道路、仓储、户籍、市场、寺院与家庭关系都可能决定命令如何抵达、战事如何延长，或接收何以在不同地区呈现不一样的速度。

本条列举的核心材料为《汉书·武帝纪》；《汉书·儒林传》；《史记·儒林列传》。这些材料较适合钉住事件的发生、官方表述或特定人物、地点，却不自动构成对全境人口、收入、兵数、信仰或动机的调查。账本的置信与材料提示是“可靠纪年；改革内容与失败过程有争议”；来源限度进一步说明“建元新政、窦太后与早期儒生政治研究”。因此，不能把条目中的政权名称、行政分类或战报修辞直接转译为固定族群和统一社会意志，也不能将制度宣布写成各地同步实施。

在叙事上，更稳妥的做法是把这一事件视为一条需要地方中介才能运作的链：中央的命令、地方官与精英的选择、运输与劳力的条件、以及普通家庭可能面对的迁徙、赋役或安全风险彼此相连。现有材料能够显示其中若干环节，却保留了大量沉默。承认这种沉默并非放弃解释，而是避免以一条编年记录覆盖不同区域和社群的差异。
